\section{The scalar certificates}\label{app:scalar}

This appendix proves the scalar assertions of Lemma~\ref{lem:scalar}.
We first verify the two-point inequality in three regions.  The
normalized checks near the two endpoints also prove concavity of $G$,
which is then used in the third region.  We next verify the polynomial
conditions that make the reflected profile reach the target and close
the dimension induction.  Finally, we prove the initial estimate on at
most three coordinates.  In the profile and base checks the noise
parameter varies over a whole interval; endpoint bounds and interval
enclosures, respectively, will control that variation.
All decimal numbers in this appendix are exact terminating decimals.
The seven choices of $(L,U,V,\alpha,\Lambda,\mu)$ and the bounds
$(r_0,r_1)$ are those of Table~\ref{ind:tab:parameters}.  A parameter
choice is fixed throughout each calculation.  At an endpoint shared by
two ranges one uses an entire row, including its profile, consistently.
The finite computations below involve only fixed polynomials and functions
of at most three real variables; there is no dimension-dependent search.
Their arithmetic recursions and all numerical input data are specified
here.  The accompanying files
\texttt{verification/interval\_certificate.cpp} and
\texttt{verification/profile\_certificate.py} implement these recursions.

\subsection{Reduction of the pair inequality}\label{sc:pair-reduction}

Write $\delta=(1-\rho)/2$, so $s^2=4\delta(1-\delta)$.
The table bounds $r_0^2\leq\delta\leq r_1^2$ follow by comparing
\[
 4r_0^2(1-r_0^2)\leq s_{\rm lo}^2,
 \qquad 4r_1^2(1-r_1^2)\geq s_{\rm hi}^2,
 \qquad r_1^2<\tfrac12.
\]
These are rational comparisons, and $a\mapsto4a(1-a)$ is increasing on
$[0,1/2]$.  We prove the pair assertion for $s>0$ and $p_0<p_1$;
when $p_0=p_1$, both its deficit and its shear are zero.

For a smooth function $F$, use divided differences
\[
 F[a,b]=\frac{F(b)-F(a)}{b-a},\qquad
 F[a,b,c]=\frac{F[b,c]-F[a,b]}{c-a},
\]
with repeated arguments interpreted by continuous extension whenever
the function is smooth there.  Put $z=p_1-p_0$, $p_t=p_0+tz$, and
\begin{equation}\label{sc:KW}
 K=-\frac{G[p_0,p_\delta,p_1]+G[p_0,p_{1-\delta},p_1]}8,
 \qquad W=1-\frac{G[p_0,p_1]}{4L}.
\end{equation}
The interpolation identity
\[
 G(p_t)-(1-t)G(p_0)-tG(p_1)
   =-t(1-t)z^2G[p_0,p_t,p_1]
\]
gives $D=s^2z^2K$.  Moreover, the shear in
Equation~\eqref{ind:eq:pair} is $zW$.  Thus the strict inequality
\begin{equation}\label{sc:Delta}
 \Delta:=K(S+\Lambda D)-\mu W^2>0
\end{equation}
implies its required nonstrict inequality.  We will also prove $D\geq0$.

The following change of variable gives smooth normalized expressions
near the endpoints, where derivatives of $G$ otherwise interfere with
interval evaluation.  For $0\leq y<1$, put
$q=\sqrt{1-y^2}$, $h=2qy$, and $e=(1+h\lambda(h))^{-1}$, and define
\begin{align}
 N(y)&=4q(y+q\lambda(h))e,\nonumber\\
 R(y)&=4q^2(q+y\lambda(h))e,\label{sc:NRM}\\
 M(y)&=4\bigl\{Ly+q[2L^2-1-2Uq^2
             +h(2LU-2Vq^2+2LVh)]\bigr\}e.\nonumber
\end{align}
Direct substitution into Equation~\eqref{ind:eq:G}, or multiplication by
$1+h\lambda(h)$, proves
\begin{equation}\label{sc:representations}
 G(y^2)=y^2N(y)=4Ly^2-y^3M(y),
 \qquad G(1-y^2)=yR(y).
\end{equation}
For every parameter row, $U<0$, $V>0$, and
$\lambda(h)\geq L+U\geq.488$ on $[0,1]$.
Consequently all the denominators just used are positive, and $N,R,M$
are smooth at the arguments used below, which stay strictly below $1$.
Also $N(y),R(y)>0$ for $0\leq y<1$.

The three expressions separate different endpoint behaviors.
$N$ removes the factor $p$ from $G(p)$ at $p=0$, while $R$ removes
the factor $\sqrt{1-p}$ at $p=1$.
The second representation of $G(y^2)$ shows that the right derivative is $G'(0)=4L$:
$M$ measures the remainder after subtracting this tangent.
Consequently the shear
$z-(G(p_1)-G(p_0))/(4L)$ vanishes for the linear function $G(p)=4Lp$.
Using $M$ in the first endpoint region preserves this cancellation
before any interval bounds are taken; using $R$ in the second removes
the square-root factor there.

\paragraph{Two endpoint regions.}
Case~1 is
\[
 p_0=v^2u^2,\quad p_1=v^2,\qquad
 (u,v)\in[0,1]\times[0,.8].
\]
Case~2 is
\[
 p_0=1-v^2,\quad p_1=1-v^2u^2,\qquad
 (u,v)\in[0,1]\times[0,.8].
\]
In both cases introduce $j\in[0,r_1]$ and set
\begin{equation}\label{sc:uwt}
 d=j(2u+(1-u)j),\quad w=u+(1-u)j,\quad
 t=\sqrt{1-(1-u)d},\quad \delta=\frac d{1+u},
\end{equation}
retaining only parameters with $r_0^2\leq\delta\leq r_1^2$.
Indeed
\[
 \delta=j^2+\frac{2u}{1+u}j(1-j),\quad
 \partial_u\delta=\frac{2j(1-j)}{(1+u)^2},\quad
 \partial_j\delta=\frac{2(u+(1-u)j)}{1+u}.
\]
It increases in both variables and, for each fixed $u$, runs continuously
from $0$ to at least $r_1^2$.  Hence every required $\delta$ is represented.
The identities
\[
 w^2=u^2+\delta(1-u^2),\qquad
 t^2=1-\delta(1-u^2)
\]
identify the two noisy mixture points (their order is reversed in
Case~2).  They also give $u\leq w,t\leq1$.

Let $m^i$, for $i=w,t$, be the divided-difference tuple of the function
$a\mapsto M(va)$ in Case~1 and $a\mapsto R(va)$ in Case~2 at the
three nodes $(u,i,1)$.  Its seven entries are indexed by
$0,1,2,01,02,12,012$; for example, $m^i_{02}$ is the first divided
difference on the endpoint nodes.  The common endpoint entries are
written $m_0,m_2,m_{02}$.  Set
\[
 \mathcal E(a,b)=\frac{a^2}{a+b}\quad(0\leq a\leq b),\qquad
 \mathcal E(0,0)=0.
\]

In Case~1 define
\begin{align}
 F&=\left(1+\frac{u^2}{1+u}\right)m_2
                  +\frac{u^3}{1+u}m_{02},\nonumber\\
 S_0&=\alpha(4L-vm_2)
               +(1-\alpha)u^2(4L-vu m_0),\nonumber\\
 P_i&=\frac{(1-\mathcal E(u,i)/(1+u))m_2
            +(i+\mathcal E(u,i))m^i_{12}
            +\mathcal E(u,i)u(m^i_{012}-m_{02}/(1+u))}{1+i},
       \quad i=w,t,\label{sc:case1-data}\\
 Z&=\frac{P_w+P_t}{8}.\nonumber
\end{align}
The finite sign check in Subsection~\ref{sc:interval} proves
\begin{equation}\label{sc:case1}
 ZS_0+4\Lambda d(1-u)(1-w^2)vZ^2
                     -\frac{\mu vF^2}{(4L)^2}>0.
\end{equation}
Here the endpoint first difference of $b^{3/2}M(v\sqrt b)$ on
$(u^2,1)$ is $F$, and its second difference on $(u^2,i^2,1)$ is $P_i$.
For clarity, the identities behind the latter statement are
\begin{align*}
 (b^{3/2})[u^2,i^2]&=i+\mathcal E(u,i),\\
 (b^{3/2})[u^2,i^2,1]&=
           \frac{1-\mathcal E(u,i)/(1+u)}{1+i},\\
 (M(v\sqrt b))[u^2,i^2,1]&=
           \frac{m^i_{012}-m_{02}/(1+u)}{(u+i)(1+i)}.
\end{align*}
They follow by factoring differences of squares and applying the
product rule for divided differences stated below.  The final formula in
Equation~\eqref{sc:case1-data} has a continuous extension at the corner $u=i=0$.
Using Equation~\eqref{sc:representations} gives, for a nondegenerate pair,
\[
 K=Z/v,\qquad W=vF/(4L),\qquad S=v^2S_0,
 \qquad s^2(1-u^2)^2=4d(1-u)(1-w^2).
\]
The left side of Equation~\eqref{sc:case1} is therefore $\Delta/v$.
Since $v>0$ for an actual pair, Equation~\eqref{sc:case1} implies
Equation~\eqref{sc:Delta} in Case~1.

In Case~2 define
\begin{align}
 F&=\frac{m_2+u m_{02}}{1+u},\nonumber\\
 Z&=\frac18\left[
 \frac{F-m^w_{12}-u m^w_{012}}{1+w}
 +\frac{u+w}{u+t}\frac{F-m^t_{12}-u m^t_{012}}{1+t}
 \right],\label{sc:case2-data}\\
 S_1&=(\alpha-\Lambda/2)u m_0
       +(1-\alpha-\Lambda/2)m_2
       +(\Lambda/2)(w m^w_1+t m^t_1).\nonumber
\end{align}
The second finite sign check is
\begin{equation}\label{sc:case2}
 ZS_1-\mu(u+w)\left(v+\frac F{4L}\right)^2>0.
\end{equation}
Applying the same difference-of-squares identities to
$\sqrt b R(v\sqrt b)$ now gives
\[
 Kv^3(u+w)=Z,\quad W=1+F/(4Lv),\quad S+\Lambda D=vS_1.
\]
Thus the left side of Equation~\eqref{sc:case2} equals $v^2(u+w)\Delta$.
For a required pair $v>0$; and $u+w=0$ would force $u=j=0$ and
$\delta=0$, excluded here.  This proves Equation~\eqref{sc:Delta} in Case~2.

\paragraph{Strict concavity, before the third region.}
The preceding sign checks also hold at repeated nodes, by their smooth
normalized formulas.  Fix $u=1$ and any permitted $j=\delta>0$. Such a value lies in the
certified interval $[0,r_1]$, since $\delta\le r_1^2<r_1$.
Then $w=t=1$.  In Case~1, $S_0=N(v)>0$ and the $Z^2$ term of
Equation~\eqref{sc:case1} vanishes, so that inequality directly forces $Z>0$.
In Case~2, $S_1=R(v)>0$, and Equation~\eqref{sc:case2} likewise directly forces
$Z>0$.  No prior sign of $K$ is used in either inference.
For $v>0$ the identities above imply
$K=-G''(p_0)/8>0$ on, respectively, $(0,.64]$ and $[.36,1)$.
These intervals cover $(0,1)$; hence $G$ is strictly concave there.
Since $G$ is continuous on $[0,1]$, its interpolation deficit is strictly
positive for any distinct endpoint pair and any interpolation weight in
$(0,1)$, including pairs containing $0$ or $1$.
It follows that $D>0$ and $K>0$ for every actual pair under consideration.
This proves the separate assertion $D\geq0$ and supplies the sign of $K$
needed next.

\paragraph{The remaining region.}
Cases~1 and~2 cover $p_1\leq.64$ and $p_0\geq.36$.
Any pair outside them is in Case~3:
\[
 p_0=l^2,\quad p_1=1-u^2,\qquad
 (l,u,r)\in[0,.6]^2\times[r_0,r_1],\qquad \delta=r^2.
\]
Put
\begin{align}
 z&=1-l^2-u^2,\quad b=r^2z,\quad
 l'=\sqrt{l^2+b},\quad u'=\sqrt{u^2+b},\nonumber\\
 P&=N(l')+\mathcal E(l,l')N[l,l'],\qquad
 Q=R(u')+uR[u,u'],\label{sc:case3-data}\\
 Z&=\frac{Q+(u+u')P}{8(1-r^2)z},\qquad
 W=1-\frac{uR(u)-l^2N(l)}{4Lz}.\nonumber
\end{align}
Here $z\geq.28$, so these denominators stay positive.  The endpoint
increments in $G$ are $bP$ and $(u'-u)Q$, respectively.  As
$b=(u+u')(u'-u)$, this proves $Z=(u+u')K$ and identifies $W$ with
Equation~\eqref{sc:KW}.  Set
\begin{align}
 Y_0={}&Z\bigl[(\alpha-\Lambda/2)uR(u)
       +(\Lambda/2)u'R(u')+(1-\alpha)l^2N(l)
       +(\Lambda/2)bP\bigr]-\mu(u+u')W^2,\nonumber\\
 Y_1={}&Z\bigl[(\alpha/2)R(u)
                   +(\Lambda/2)(u'-u)P\bigr]-\mu W^2,\label{sc:Y}\\
 Y_2={}&Z\bigl[(\Lambda/2)Q
                   +(\Lambda/2)(u'-u)P\bigr]-\mu W^2.\nonumber
\end{align}
The third finite check proves the disjunction
\begin{equation}\label{sc:case3}
 Y_0>0\quad\hbox{or}\quad (Y_1>0\ \hbox{and}\ Y_2>0).
\end{equation}
The first alternative suffices because $Y_0=(u+u')\Delta$.
In the second alternative take the convex combination of the brackets
in $Y_1,Y_2$ with weights $2u/(u+u')$ and $(u'-u)/(u+u')$.
It is
\[
 \frac{\alpha uR(u)+(\Lambda/2)(u'-u)Q+(\Lambda/2)bP}{u+u'}
 =\frac{S+\Lambda D-(1-\alpha)l^2N(l)}{u+u'}
 \leq\frac{S+\Lambda D}{u+u'}.
\]
For $s>0$ we have $u+u'>0$, and the concavity already proved gives
$Z=(u+u')K>0$.  Therefore the second alternative also implies
Equation~\eqref{sc:Delta}.  The three regions prove Equation~\eqref{ind:eq:pair}, subject
only to the following explicit finite sign checks.

\subsection{Interval arithmetic for the three pair regions}\label{sc:interval}

\paragraph{Enclosures and derivatives.}
All intervals may have rational endpoints.  The particular certificate
rounds outwards after every operation to multiples of $Q_*^{-1}$,
where $Q_*=10^{12}$.  Store an interval as $[a/Q_*,b/Q_*]$ with integers
$a\leq b$.  Addition and subtraction use their ordinary endpoint rules.
For multiplication, take the least and greatest of the four endpoint
products and round their quotients by $Q_*$ down and up.  For a positive
interval its reciprocal endpoints are
\[
 Q_*^{-1}\left\lfloor\frac{Q_*^2}{b}\right\rfloor,
 \qquad Q_*^{-1}\left\lceil\frac{Q_*^2}{a}\right\rceil.
\]
For a nonnegative interval the square-root endpoints are
$Q_*^{-1}\lfloor\sqrt{aQ_*}\rfloor$ and
$Q_*^{-1}\lceil\sqrt{bQ_*}\rceil$; integer square comparison computes
these exactly.  Failed positivity or domain conditions never certify
a box.  A power with a nonnegative small integer exponent is evaluated
by repeated interval multiplication, including for squares; this may
overestimate but remains an enclosure.

An interval expression may be augmented by its three first derivatives.
Use the sum and product rules, and
\[
 \partial_i(X^{-1})=-X^{-2}\partial_iX,\qquad
 \partial_i\sqrt X=\frac{\partial_iX}{2\sqrt X}.
\]
The latter derivative is used only when the lower bound for the root is
strictly positive.  In Case~3 the radicands defining $l',u'$ are known
nonnegative on the whole box; direct \emph{value} evaluation may therefore
replace a negative interval lower endpoint by zero.  Derivative evaluation
does not use this clipping and fails at an uncertified root singularity.

The function $\mathcal E(a,b)$ has a special rule on $0\leq a\leq b$.
Enclose $h=a/(a+b)$ by $[0,1/2]$, intersecting with direct interval
division if $a+b$ has a positive lower bound.  Its value is enclosed by
$ah$, and its differential by
\begin{equation}\label{sc:E-derivative}
 d\mathcal E=h(2-h)\,da-h^2\,db.
\end{equation}
At $(0,0)$ the function is continuous and Lipschitz on this cone, so the
same difference bounds hold by limits from its interior; differentiability
at that point is not assumed.  The required order relations $u\leq w,t$
and $l\leq l'$ hold throughout each parameter box, not just on its
$\delta$-restricted subset.  Thus these enclosures apply on the line
segments used in the mean-value estimate below.

\paragraph{Tuples without division by node distances.}
For three nodes let $a$ denote the tuple of values and divided differences
of a function.  Addition is componentwise.  For $c=ab$, the product rule is
\begin{equation}\label{sc:tuple-product}
 c_p=a_pb_p,\quad
 c_{pq}=a_pb_{pq}+a_{pq}b_q\ (p<q),\quad
 c_{012}=a_0b_{012}+a_{01}b_{12}+a_{012}b_2.
\end{equation}
For $a=c^{-1}$ use
\begin{equation}\label{sc:tuple-reciprocal}
 a_p=1/c_p,\quad a_{pq}=-c_{pq}a_pa_q,\quad
 a_{012}=-a_0(c_{012}a_2+c_{01}a_{12});
\end{equation}
for $a=\sqrt c$ use
\begin{equation}\label{sc:tuple-root}
 a_p=\sqrt{c_p},\quad a_{pq}=\frac{c_{pq}}{a_p+a_q},\quad
 a_{012}=\frac{c_{012}-a_{01}a_{12}}{a_0+a_2}.
\end{equation}
The last two formulas solve $ac=1$ and $a^2=c$ using
Equation~\eqref{sc:tuple-product}.  They remain valid at repeated nodes by
continuity and use no inverse node distance.  All entries can themselves
be intervals, or intervals augmented by the three derivatives.  A constant
tuple has that constant in its three value entries and zero in the other
entries.  Write $X(a,b,c,g)$ for the tuple with value entries $(a,b,c)$,
all three first differences $g$, and second difference zero.

To specify the evaluation completely, the tuple-valued function
$\operatorname{table}(x,k)$ is obtained by the following straight-line
recipe; $\operatorname{sq}$ means multiplication of an expression by itself.
\begin{verbatim}
q2 = 1-x*x; q = sqrt(q2); h = 2*(q*x)
a = L+h*(U+V*h); e = inv(1+h*a)
if k==1:
  return 4*(L*x+q*((2*sq(L)-1)-2*U*q2
                  +h*((2*L*U)-2*V*q2+(2*L*V)*h)))*e
if k==2: return 4*(q2*(q+x*a))*e
return 4*(q*(x+q*a))*e
\end{verbatim}
These are exactly $M,R,N$ in that order.  The expressions to test are
specified below, retaining the evaluation order of the certificate.
Tuple subscripts are strings of node indices, so $12$ and $012$ denote
first and second divided differences.  Products and quotients on one
line associate from left to right.  The symbols \texttt{alpha},
\texttt{Lambda}, \texttt{mu}, and \texttt{E} stand for
$\alpha,\Lambda,\mu,\mathcal E$, respectively.
\begin{verbatim}
term(m,u,i):
  e = E(u,i); o = inv(1+u)
  return ((1-o*e)*m[2]+(i+e)*m[12]
                         +e*u*(m[012]-o*m[02]))/(1+i)

\end{verbatim}
\begin{samepage}
\begin{verbatim}
expr(u,v,j,k):
  d = j*(2*u+(1-u)*j); w = u+(1-u)*j
  t = sqrt(1-(1-u)*d)
  m = table(X(v*u,v*w,v,v),k)
  n = table(X(v*u,v*t,v,v),k)
  if k==1:
    F = (1+u*u/(1+u))*m[2]+u*u*u*m[02]/(1+u)
    Z = .125*(term(m,u,w)+term(n,u,t))
    H0 = alpha*(4*L-v*m[2])
                            +(1-alpha)*u*u*(4*L-v*u*m[0])
    di = d*(1-u); wi = 1-sq(w)
    return [Z*H0+(4*Lambda)*di*wi*v*sq(Z)
                                      -mu*v*sq(F)/sq(4*L)]
  half = .5*Lambda
  F = (m[2]+u*m[02])/(1+u)
  Z = .125*((F-m[12]-u*m[012])/(1+w)
                 +(u+w)/(u+t)*(F-n[12]-u*n[012])/(1+t))
  H0 = (alpha-half)*u*m[0]+(1-alpha-half)*m[2]
                                       +half*(w*m[1]+t*n[1])
  return [Z*H0-mu*(u+w)*sq(v+F/(4*L))]

\end{verbatim}
\end{samepage}
\par\pagebreak[0]
\begin{samepage}
\begin{verbatim}
expr3(l,u,r):
  z = 1-sq(l)-sq(u); b = sq(r)*z
  lp = nroot(sq(l)+b); up = nroot(sq(u)+b)
  n = table(X(l,l,lp,1),3); a = table(X(u,u,up,1),2)
  P = n[2]+E(l,lp)*n[02]; Q = a[2]+u*a[02]
  Z = (Q+(u+up)*P)/(8*(1-sq(r))*z)
  W2 = sq(1-(u*a[0]-l*l*n[0])/(4*L*z))
  half = .5*Lambda
  Y0 = Z*((alpha-half)*u*a[0]+half*up*a[2]
                    +(1-alpha)*l*l*n[0]+half*b*P)
                                            -mu*(u+up)*W2
  rest = half*(up-u)*P
  Y1 = Z*((.5*alpha)*a[0]+rest)-mu*W2
  Y2 = Z*(half*Q+rest)-mu*W2
  return [Y0,Y1,Y2]
\end{verbatim}
\end{samepage}
Here \texttt{nroot} is the square-root operation with the value-only
clipping already specified; on derivative tuples it is the ordinary
strict-domain root.  These recipes are the normalized expressions in
Equations~\eqref{sc:case1}, \eqref{sc:case2}, and \eqref{sc:Y}.

\paragraph{The box test.}
An initial box has coordinates $(u,v,j)$ in Cases~1 and~2 and
$(l,u,r)$ in Case~3, with the ranges specified above.
A subdivision bisects all three coordinates, giving eight children.
First exclude a box in Cases~1 and~2 if the interval between the low-low
and high-high corner values of $\delta(u,j)$ misses $[r_0^2,r_1^2]$.
The monotonicity in Equation~\eqref{sc:uwt} proves this exclusion sound, with the
corner arithmetic rounded outwards.

For every remaining scalar expression $F$ one may certify $F>0$ either
by a positive lower endpoint of its interval value on the whole box, or
by the following derivative test.  Let $c_i,h_i$ be the exact center and
half-width of its coordinate intervals, and normalize the coordinates by
$x_i=c_i+h_i\xi_i$ with $|\xi_i|\leq1$.  Initialize the derivative of
coordinate $x_i$ with respect to $\xi_i$ to an outward enclosure of $h_i$,
and its other derivatives to zero.  If $I_c$ encloses $F(c)$ and $J_i$
encloses $\partial F/\partial\xi_i$ on the full box, the sufficient test is
\begin{equation}\label{sc:center-test}
 \inf I_c>\sum_{i=1}^3\max(|\inf J_i|,|\sup J_i|).
\end{equation}
The mean-value theorem, with the limiting argument for
Equation~\eqref{sc:E-derivative} when needed, proves the test.
In Case~3 different components can use different tests: the box is good
if the first component is positive, or if both the second and third are
positive.  It is unnecessary for the same branch of the disjunction to
hold on different boxes.

More formally, with $d$ the remaining permitted depth, the entire
finite decision procedure is
\begin{verbatim}
verify(box,d):
  if the box is excluded by the delta restriction: return true
  attempt interval values and, if needed, center/derivative bounds
  if the Case 1 or 2 expression is certified positive: return true
  if Case 3 has component 0 positive, or components 1 and 2 positive:
      return true
  # An invalid arithmetic operation never certifies a box.
  if d==0: return false
  return AND(verify(child,d-1) for the eight children)
\end{verbatim}
In the supplied implementation any invalid attempt immediately proceeds
to subdivision; it does not turn a failed bound into a certificate.
A sign-certified leaf proves the required inequality on its entire box;
an excluded leaf has empty intersection with the required delta range.
Induction up the finite subdivision tree therefore proves the inequalities
on the initial region subject to that restriction.

\begin{table}[htbp]
\centering
\begin{tabular}{lrrrrrrr}
\toprule
Parameter group &0&1&2&3&4&5&6\\
\midrule
Case 1 &8&8&8&7&8&8&8\\
Case 2 &8&8&8&7&8&8&8\\
Case 3 &14&7&7&7&7&7&7\\
\bottomrule
\end{tabular}
\caption{Sufficient maximum subdivision depths for the pair certificates;
the initial box has depth zero.}\label{sc:tab:interval-depths}
\end{table}

Evaluation of the recursions with $Q_*=10^{12}$ gives \texttt{true}
for all 21 initial boxes at the depths in
Table~\ref{sc:tab:interval-depths}.  This is a finite list of integer
comparisons, as specified above.  In particular it proves the strict
sign assertions in Equations~\eqref{sc:case1}, \eqref{sc:case2}, and \eqref{sc:case3},
including the normalized boundary values used for the concavity argument.
The recorded execution of the accompanying C++ certificate has no failed
terminal leaf.

\paragraph{Arithmetic range and trust.}
The mathematical procedure can use arbitrary-precision integers.  The
supplied C++ implementation instead guards each stored scaled endpoint
by $|a|,|b|\leq10^{18}$.  Endpoint sums and differences before the guard
are at most $2\cdot10^{18}$ in magnitude, and the sum of three derivative
magnitudes in Equation~\eqref{sc:center-test} is at most $3\cdot10^{18}$; these
fit signed 64-bit integers.  Each endpoint product is at most $10^{36}$,
below the signed 128-bit maximum; the reciprocal numerator is $10^{24}$,
and root radicands are at most $10^{30}$.  The integer root search uses
upper bound $10^{15}+1$, whose square also fits signed 128 bits.
At depths at most $14$, dyadic grid indices are below $2^{14}$; the
pre-conversion grid products in the source are below $2\cdot10^{16}$.
Its unguarded integer constructor receives only the displayed small
constants.  Thus intermediate operations as well as stored values fit
their integer types.  An overflow guard failure causes subdivision,
not acceptance.  The computational verification uses ordinary exact
integer arithmetic and reviewed source execution; it does not require
floating-point rounding assumptions or claim proof-assistant verification.

\subsection{The seventeen polynomial profiles}\label{sc:profiles}

It remains to choose a profile that satisfies the reflected comparison
and the induction closure condition, Lemma~\ref{lem:scalar}(2)--(3).
We use seventeen polynomials of the form
\begin{equation}\label{sc:profile-def}
 B(m)=(1-m^2)P(m),\qquad
 P(m)=1+10^{-8}\sum_{j\geq1}c_j\bigl(T_j(m)-T_j(0)\bigr),
\end{equation}
where $T_0=1$, $T_1=m$, and $T_{j+1}=2mT_j-T_{j-1}$.
The form in Equation~\eqref{sc:profile-def} gives $B(0)=1$ and $B(\pm1)=0$
exactly.  The coefficient data below assign each polynomial to a closed
interval $[s_0,s_1]$ contained in one parameter group.  Before giving
those data, we reduce the remaining properties to five sign checks.

\paragraph{The required signs and their implications.}
Use the polynomial quotients
\begin{align}
 E_B(m)&=\frac{(B(m)+B(-m))/2-1}{m^2},\nonumber\\
 C(m,x)&=\frac{B(m+x)-B(m-x)}{2x},\qquad
 A(m,x)=\frac{2B(m)-B(m+x)-B(m-x)}{2x^2},\label{sc:quotients}\\
 J_s&=1-\frac{s}{2L}C,\qquad
 \mathcal V_s=\mu J_s^2
       -A\bigl[B(m)+(2\alpha-1)xC+(\Lambda-1)x^2A\bigr].\nonumber
\end{align}
They are polynomial identities, including at $m=0$ or $x=0$.
For $B(m)=\sum b_km^k$, explicit coefficient recipes are
\begin{align*}
 C&=\sum_k\sum_{\substack{1\leq j\leq k\\j\ {\rm odd}}}
       \binom{k}{j}b_km^{k-j}x^{j-1},\\
 A&=-\sum_k\sum_{\substack{2\leq j\leq k\\j\ {\rm even}}}
       \binom{k}{j}b_km^{k-j}x^{j-2},\qquad
 E_B=\sum_{\substack{k\geq2\\k\ {\rm even}}}b_km^{k-2}.
\end{align*}
There is thus no evaluation of a quotient by a vanishing variable.

For each of the seventeen profiles on its assigned interval, we verify
\begin{align}
 P(m)&>0 &&(-1\leq m\leq1),\nonumber\\
 1+s_1E_B(m)(1+H(m))&>0 &&(0\leq m\leq1),\label{sc:profile-signs}\\
 J_{s_1},\ \mathcal V_{s_0},\ \mathcal V_{s_1}&>0
       &&(-1\leq m\leq1,\ 0\leq x\leq b(m)).\nonumber
\end{align}
Here the piecewise affine upper boundary $b(m)$ is completely specified
by the successive knot pairs, multiplied by $16$,
\begin{equation}\label{sc:knots}
 (-16,0),\ (-14,2),\ (-6,6),\ (6,6),\ (14,2),\ (16,0).
\end{equation}
In particular $b(m)\leq1-|m|$, so all arguments $m\pm x$ are in $[-1,1]$.

These sign assertions suffice for the reflection and profile parts of
the scalar lemma, as follows.  First $B=(1-m^2)P\geq0$.
The polynomial $E_B$ is even, and $1-H(m)=m^2/(1+H(m))$, whence
\begin{equation}\label{sc:reflection-difference}
 \frac{s}{2}\bigl(B(m)+B(-m)\bigr)-(H(m)-1+s)
   =m^2\left(sE_B(m)+\frac1{1+H(m)}\right).
\end{equation}
For $E_B<0$ the certified bound at $s_1$ implies the one at any
$0\leq s\leq s_1$; for $E_B\geq0$ it is immediate.  Thus
Equation~\eqref{ind:eq:reflection-profile} follows, including $m=0$, where both
sides equal $s$, and $m=\pm1$.

For fixed $(m,x)$, $J_s$ is affine in $s$, $J_0=1$, and the certificate
gives $J_{s_1}>0$.  Therefore $J_s>0$ on $[0,s_1]$, and its square is
monotone there: decreasing when $C>0$, increasing when $C<0$, constant
when $C=0$.  Consequently
\[
 J_s^2\geq\min(J_{s_0}^2,J_{s_1}^2)
       \qquad(s_0\leq s\leq s_1).
\]
The expression subtracted from $\mu J_s^2$ in $\mathcal V_s$ is independent
of $s$ within the row.  The two endpoint certificates therefore imply
$\mathcal V_s>0$ throughout the row.  This proves
Equation~\eqref{ind:eq:profile} for either sign of $A$ or $C$.
It is the positivity of the affine factor that justifies this step;
endpoint positivity of a general quadratic alone would not suffice.

\subsection{Profile data and exact sign verification}\label{sc:bernstein}

\paragraph{Coefficient data.}
Each line below lists $1000s_1$ followed by $(c_1,c_2,\ldots)$;
unlisted coefficients are zero.  Its interval is $[s_0,s_1]$, where
$s_0$ is the previous endpoint, starting at zero.  The group is the
first row of Table~\ref{ind:tab:parameters} whose upper endpoint is at
least $s_1$; no profile interval straddles a group boundary.

\begingroup\small
\noindent\textbf{5:}
26797227, 21421001, 5128970, 9234452, 692060, 4904060, -868991,
3430822, -1620202, 2642536, -1572003, 1611382, -790463, 573913,
-196263, 125941, -28650, 17140.\par\smallskip
\noindent\textbf{17:}
26545421, 20782313, 4702871, 8733755, 262302, 4632398, -1223407,
3219225, -1751136, 2367765, -1499679, 1323103, -650254, 364146,
-98062, 32125.\par\smallskip
\noindent\textbf{39:}
23422612, 23811053, 2720392, 9481334, 342747, 3521407, 293182,
1243280, 94220, 628296, -211618, 367918, -115140, 71341.\par\smallskip
\noindent\textbf{70:}
28627418, 18132350, 6560313, 4688850, 3090858, 280132, 2014606,
-721794, 1092862, -413975, 251954, -30987.\par\smallskip
\noindent\textbf{108:}
25845226, 21339334, 2692615, 7536216, -587180, 2879744, -642549,
859994, -135519, 99765.\par\smallskip
\noindent\textbf{160:}
28111431, 22222528, 3686628, 6969749, 260698, 2311757, -230942,
626183, -39197, 65219.\par\smallskip
\noindent\textbf{220:}
31980620, 22310973, 3943312, 6252289, 221019, 1801879, -204196,
345834.\par\smallskip
\noindent\textbf{300:}
37050025, 22730221, 6023971, 4990611, 1410047, 1047683, 133030,
195650.\par\smallskip
\noindent\textbf{400:}
46891988, 25975928, 8279226, 5707508, 1896038, 1189721, 225195,
190389.\par\smallskip
\noindent\textbf{490:}
57774207, 30017282, 10752796, 6552311, 2404507, 1338857, 310628,
193643.\par\smallskip
\noindent\textbf{570:}
67261317, 34557385, 13515609, 7744945, 3024441, 1590294, 428436,
214704.\par\smallskip
\noindent\textbf{630:}
76105766, 38746702, 15974003, 8801748, 3567377, 1794423, 517893,
232715.\par\smallskip
\noindent\textbf{700:}
92621845, 50731247, 23924150, 14690765, 6918336, 4309886, 1672440,
1096916, 239750, 209175.\par\smallskip
\noindent\textbf{750:}
100973193, 55145076, 26634158, 15924290, 7601086, 4546032, 1797080,
1099352, 246042, 192495.\par\smallskip
\noindent\textbf{800:}
114390772, 63949520, 33959789, 20265973, 11129205, 6354835, 3302519,
1693517, 778655, 305292, 118236.\par\smallskip
\noindent\textbf{827:}
126357200, 73807375, 40646383, 25693312, 14417512, 9123323, 4734690,
2955541, 1274070, 786030, 225052, 144665.\par\smallskip
\noindent\textbf{840:}
140455073, 85610577, 50361189, 33435466, 20391097, 13705759, 8022488,
5381807, 2815510, 1892888, 793925, 548406, 137909, 116217.\par
\endgroup

We now verify the five sign conditions in
Equation~\eqref{sc:profile-signs} for these data.

\paragraph{Bernstein bounds.}
For a polynomial in coordinates $(u,v)\in[0,1]^2$, of degree at most
$N$ in $u$ and $K$ in $v$, write $R(u,v)=\sum r_{ij}u^iv^j$.
The formula of Appendix~\ref{app:arithmetic} gives its Bernstein
coefficients of orders $(N,K)$ as
\begin{equation}\label{sc:bernstein-transform}
 \beta_{nk}=\sum_{i=0}^n\sum_{j=0}^k
 r_{ij}\frac{\binom ni}{\binom Ni}\frac{\binom kj}{\binom Kj},
 \qquad 0\leq n\leq N,\quad0\leq k\leq K.
\end{equation}
The Bernstein enclosure proved in Appendix~\ref{app:arithmetic} puts
the polynomial between the least and greatest coefficients on the full box.
For a univariate polynomial set $K=0$ and omit the second coordinate.

For $P$ start with $[-1,1]$; for $E_B$ use $[0,1]$.
For the other three polynomials use each consecutive pair of knots
$(a,c),(a+h,c+d)$ from Equation~\eqref{sc:knots}, in their actual rational units,
and substitute
\[
 m=a+hu,\qquad x=v(c+du),\qquad (u,v)\in[0,1]^2.
\]
For each polynomial $S$ to be checked, sufficient orders are
$N=\deg_{\mathrm{tot}}S$ and $K=\deg_xS$ before substitution,
because $m^ix^j$ has $u$-degree at most $i+j$
and $v$-degree $j$.  These rectangles cover the entire closed domain.
Subdivision bisects both coordinates.  Equivalently, on a subdivided
rectangle $m\in[a,a+h]$, $t\in[c,c+d]$, substitute
\[
 m=a+hu,\quad t=c+dv,\quad
 x=t\bigl(b(a)+(b(a+h)-b(a))u\bigr).
\]
All these substitutions and Equation~\eqref{sc:bernstein-transform} are rational.

For the columns $P,J,\mathcal V$ every final coefficient must be strictly
positive.  For the $E_B$ column, on each interval with lower endpoint $a$,
each negative coefficient $e$ must instead satisfy
\begin{equation}\label{sc:E-test}
 1+s_1e>0,\qquad (1+s_1e)^2>(s_1e)^2(1-a^2).
\end{equation}
These two comparisons prove $1+s_1e(1+H(a))>0$ with the sign checked
before squaring.  Since $H(m)\leq H(a)$ on that interval and the polynomial
is bounded below by its least Bernstein coefficient, this proves the
second line of Equation~\eqref{sc:profile-signs} everywhere on the interval.
If its least coefficient is nonnegative there is nothing to check.

\begin{table}[htbp]
\centering
\begin{tabular}{rrrrrrr}
\toprule
&\multicolumn{5}{c}{Bisection depths}&Base lower bound\\
\cmidrule(lr){2-6}\cmidrule(l){7-7}
$1000s_1$&$P$&$E_B$&$J_{s_1}$&$\mathcal V_{s_0}$&$\mathcal V_{s_1}$&\\
\midrule
5&1&0&0&3&4&7\\
17&1&0&0&2&3&7\\
39&1&0&0&1&1&8\\
70&1&0&0&1&2&10\\
108&0&0&0&1&1&7\\
160&0&0&0&0&1&14\\
220&0&0&0&0&0&78\\
300&0&0&0&0&0&69\\
400&0&0&0&0&0&146\\
490&0&0&0&0&0&224\\
570&0&2&0&0&0&233\\
630&0&3&0&0&0&255\\
700&1&4&0&0&0&236\\
750&1&5&0&1&0&222\\
800&1&5&0&1&0&191\\
827&1&6&0&1&1&177\\
840&1&6&0&1&1&176\\
\bottomrule
\end{tabular}
\caption{Sufficient bisection depths for the five profile sign checks,
and lower bounds for the scaled base-case difference defined in
Equation~\eqref{sc:base-column}.}\label{sc:tab:profile-depths}
\end{table}

For each row in Table~\ref{sc:tab:profile-depths}, perform the five
Bernstein checks with its listed depths and coefficients.  The rational
comparisons in Equations~\eqref{sc:bernstein-transform}--\eqref{sc:E-test} all hold.
One can stop early on boxes whose signs are already certified: subsequent
Bernstein coefficients are convex combinations of the previous ones.
This supplies a finite arithmetic verification of Equation~\eqref{sc:profile-signs}.
For completeness, the following integer scalings specify exactly how the
accompanying Python certificate performs these rational calculations.

\paragraph{Exact integer scaling and subdivision.}
Put $D_*=10^8$, and write
\[
 \ell=1000L,\quad a_*=1000\alpha,\quad
 \lambda_*=1000\Lambda,\quad \nu=1000\mu,\quad S_i=1000s_i.
\]
The Chebyshev recurrence first forms the integer polynomials
$\widehat P=D_*P$, $\widehat B=D_*B$.
Applying the coefficient recipes following Equation~\eqref{sc:quotients} to $\widehat B$
gives $\widehat E=D_*E_B$, $\widehat A=D_*A$, $\widehat C=D_*C$.
The five polynomials supplied to the coefficient checker are
\begin{align*}
 &\widehat P,\quad\widehat E,\quad\widehat J(S_1),
       \quad\widehat V(S_0),\quad\widehat V(S_1),\\
 \widehat J(S)&=2\ell D_*-S\widehat C,\\
 \widehat V(S)&=\nu\widehat J(S)^2-(2\ell)^2\widehat A
       \bigl[1000\widehat B+(2a_*-1000)x\widehat C
                   +(\lambda_*-1000)x^2\widehat A\bigr].
\end{align*}
Thus their positive scale factors relative to the mathematical
polynomials $P,E_B,J_s,\mathcal V_s,\mathcal V_s$ are, respectively,
\begin{equation}\label{sc:scalings}
 D_*,\quad D_*,\quad 2\ell D_*,\quad
 1000(2\ell)^2D_*^2,\quad1000(2\ell)^2D_*^2.
\end{equation}

For an integer polynomial $p(m,x)=\sum p_{ij}m^ix^j$ and an initial map
\[
 m=(a+bu)/d_*,\qquad x=v(c+du)/d_*,
\]
multiply the power coefficients after substitution by $d_*^N$.
Explicitly add
\[
 p_{ij}d_*^{N-i-j}\binom{i}{k}a^{i-k}b^k
                         \binom{j}{l}c^{j-l}d^l
\]
to the coefficient of $u^{k+l}v^j$, for $0\leq k\leq i$,
$0\leq l\leq j$.  Transform the first axis by
\[
 T^{(N)}_{ki}=\binom{k}{i}\frac{N!}{\binom Ni}
\]
and the second by $T^{(K)}$.  The divisions are exact, since
$N!/\binom Ni=i!(N-i)!$.  The resulting integer matrix is
$d_*^NN!K!$ times its Bernstein coefficient array, a positive scale.

For an axis of order $n$, its left-half subdivision matrix has entries
\[
 L^{(n)}_{ik}=2^{n-i}\binom{i}{k}\qquad(0\leq i,k\leq n),
\]
with $\binom{i}{k}=0$ for $k>i$.  This is $2^n$ times the usual
de Casteljau half-subdivision.  Reversing both row and column indices
gives the right-half matrix.  Apply both matrices on both axes at each
subdivision, retaining their positive common factors.  Order-zero axes
have only one child.  Positivity tests are therefore tests of integers.

For the $E_B$ column with depth $d$, set $q_*=2^d$ and, on
$[i/q_*,(i+1)/q_*]$, let $c$ be the least integer coefficient produced
by the preceding transformation with $d_*=q_*$.  Its scale relative to
$E_B$ is $D_*q_*^NN!$.  Define
\[
 h_*=-S_1c,\qquad C_*=1000D_*q_*^NN!.
\]
If $h_*\leq0$, the test passes immediately.  Otherwise the exact tests are
\begin{equation}\label{sc:E-integer}
 C_*>h_*,\qquad
 ((C_*-h_*)q_*)^2>h_*^2(q_*^2-i^2).
\end{equation}
They are precisely Equation~\eqref{sc:E-test} for the minimum coefficient.
This accounts for every factor of $1000$, $2$, and $D_*$ in the
certificate.  Python uses unbounded integers and exact rational
evaluation here, so integer overflow is not an additional hypothesis.

\subsection{The three-coordinate base case}\label{sc:base}

We finish by proving Lemma~\ref{lem:scalar}(4).  An increasing
function on fewer than three coordinates can be viewed as a function on
three.  Constants satisfy Equation~\eqref{ind:eq:base} because
$G(0)=G(1)=B(\pm1)=0$.  Consider a nonconstant function whose positive
set has size at most four.  Up to coordinate permutations, a positive
set of size one is the top vertex; one of size two adds one of its three
neighbors; and one of size three adds two of those neighbors.  A vertex
with only one positive coordinate would force its entire upper subcube
of size four to be positive.  Hence a positive set of size four either
is that subcube, giving a dictator, or consists of the three vertices
with two positive coordinates and the top vertex, giving majority.
Larger positive sets are obtained from these by the duality
$f^*(x)=-f(-x)$.

For a dictator, $T_\rho f=\rho x_i$ and
Equation~\eqref{ind:eq:G-reflection} gives equality in
Equation~\eqref{ind:eq:base}, since $B(0)=1$.  For majority,
\[
 f(x)=\frac{x_1+x_2+x_3-x_1x_2x_3}{2}.
\]
Again the asymmetric term averages to zero by input reversal.  Evaluating
the two absolute-value classes of $T_\rho f$ gives
\[
 \frac{\E H(T_\rho f)}s
 =\frac{s\sqrt{3+s^2}+3\sqrt{8-5s^2+s^4}}8\ge1.
\]
Indeed, $s\sqrt{3+s^2}\ge2s^2$ for $0\le s\le1$, while, with $u=s^2$,
\[
 (8-5u+u^2)-\left(2+\frac34(1-u)\right)^2
 =\frac7{16}(1-u)^2\ge0.
\]
The resulting lower bound for the numerator is
$33/4-s^2/4\ge8$.  Thus only the three unbalanced positive sets
$\{7\}$, $\{7,6\}$, and $\{7,6,5\}$ in binary coding (identify $-1$ with
$0$ and $+1$ with $1$, then read the three bits as a binary integer), and their duals,
require the scalar base certificates below.
In particular the base assertion treats both orientations of $B$;
no symmetry of $B$ is assumed.

For $D\in\{\{7\},\{7,6\},\{7,6,5\}\}$ put
$m_D=|D|/4-1$ and $\sigma_i=2\mathbf1_D(i)-1$, where
$i\in\{0,\ldots,7\}$.  If $d_{ij}$ denotes Hamming distance, define
\begin{align}
 \delta&=\frac{s^2}{2(1+\sqrt{1-s^2})},\qquad
 I_i=\frac14\sum_{\mathbf1_D(i)\ne\mathbf1_D(j)}
            \delta^{d_{ij}-1}(1-\delta)^{2-d_{ij}},\nonumber\\
 e_i&=s^2I_i,\qquad h_i=2\sqrt{I_i(1-e_i)},\label{sc:base-formulas}\\
 q_i&=\frac{\sigma_i(1-2e_i)}
                {1+sh_i\bigl(L+sh_i(U+sh_iV)\bigr)}.\nonumber
\end{align}
A noise transition from $i$ to an opposite point $j$ has probability
$\delta^{d_{ij}}(1-\delta)^{3-d_{ij}}$.  Since
$s^2=4\delta(1-\delta)$, $e_i$ is exactly the boundary-crossing probability.
Thus $d_i:=T_\rho f(i)=\sigma_i(1-2e_i)$,
$H(d_i)/s=h_i$, and
$d_i/(1+H(d_i)\lambda(H(d_i)))=q_i$.
It follows that the two quantities to bound below are exactly
\begin{equation}\label{sc:base-normalized}
 \frac18\sum_{i=0}^7h_i(1+\eta q_i)-B(\eta m_D),
                   \qquad \eta\in\{-1,1\}.
\end{equation}
For the dual, global input reversal merely reindexes the average and
negates its noised output, which explains the sign $\eta$ in both places.

There is no division by $s$ in Equation~\eqref{sc:base-formulas}.  Opposite points
are distinct, so $1\leq d_{ij}\leq3$; the sole possible negative power is
$(1-\delta)^{-1}$, nonsingular over all the row intervals.
These formulas have continuous extensions to $s=0$, even if $I_i=0$.
In that limit $I_i$ is one quarter of the number of opposite neighbors,
$e_i=0$, $h_i$ is the square root of that number, and $q_i=\sigma_i$.
The two limits of $\E G/s$ are, for the three listed sets in order,
\[
 (\sqrt3/4,\ 3/4),\quad
 (\sqrt2/2,\ 1),\quad
 ((1+2\sqrt2)/4,\ (3+\sqrt2)/4).
\]
Accordingly the first closed row interval is a valid domain for the
normalized certificate, although Equation~\eqref{ind:eq:base} itself only needs
$s>0$.

For each profile row divide $[s_0,s_1]$ into exactly $128$ equal
subintervals.  On each, evaluate Equation~\eqref{sc:base-formulas} in the displayed
order with the interval arithmetic at scale $Q_*=10^{12}$ from
Subsection~\ref{sc:interval}, and form the average in
Equation~\eqref{sc:base-normalized}.  All signed profile evaluations
$B(\eta m_D)$ are exact rational numbers.  The base column of
Table~\ref{sc:tab:profile-depths} is a lower bound, simultaneously for
all $128$ intervals, all three sets, and both signs, for
\begin{equation}\label{sc:base-column}
 \left\lfloor10^4\left(
 \inf\left[\frac18\sum_{i=0}^7h_i(1+\eta q_i)\right]
                              -B(\eta m_D)\right)\right\rfloor.
\end{equation}
Equivalently, compare $10^4(\text{interval lower bound}-B(\eta m_D))$
directly with the listed integer.  No approximate floor operation is
needed.  The $128$ outward-enclosed intervals cover the entire row,
including both endpoints.  Every entry in the base column is positive,
which proves Equation~\eqref{ind:eq:base} for the unbalanced cases and hence for
all increasing functions on at most three coordinates.

Together the pair-region certificates, the profile signs, and the base
certificate prove all assertions of Lemma~\ref{lem:scalar}.
